\sectionkicker{Source-backed technical notes}
\subsection*{Coding ModelからAgentic Engineeringへ: Technical Notes}
本欄は記事本文で圧縮した一次資料上の情報を、比較・再検証しやすい形で整理したものである。時系列、確認済みの主張、留意点、一次資料URLを掲載する。
\medskip
\subsection*{Theme at a glance}
\addcontentsline{toc}{subsection}{Theme at a glance}
\begin{center}
\footnotesize
\begin{tabularx}{\linewidth}{@{}p{0.20\linewidth}p{0.10\linewidth}p{0.14\linewidth}X@{}}
\toprule
資料 & 位置づけ & 種別 & 時系列 \\
\midrule
OpenAI Codex February stack: app, GPT-5.3-Codex, System Card, and Codex-Spark & 主要資料 & モデル更新 & 2026-02-02 (Agent製品公開); 2026-02-05 (モデル公開); 2026-02-05 (System Card公開); 2026-02-12 (研究Preview) \\
Claude Opus 4.6 / Sonnet 4.6 & 主要資料 & モデル更新 & 2026-02-05 (モデル公開); 2026-02-17 (モデル公開) \\
\bottomrule
\end{tabularx}
\end{center}
\normalsize

\begin{technicalnote}{OpenAI Codex February stack: app, GPT-5.3-Codex, System Card, and Codex-Spark}{主要資料}
\begingroup
% reader-facing Technical Notes break policy
\widowpenalty=10000
\clubpenalty=10000
\displaywidowpenalty=10000
\begin{tabularx}{\linewidth}{@{}>{\bfseries}p{0.22\linewidth}X@{}}
組織 & OpenAI \\
種別 & モデル更新 \\
時系列 & 2026-02-02 (Agent製品公開); 2026-02-05 (モデル公開); 2026-02-05 (System Card公開); 2026-02-12 (研究Preview) \\
\end{tabularx}
\smallskip
{\bfseries 一次資料から整理したtechnical points}
\begin{itemize}[leftmargin=1.5em,itemsep=0.35em]
\item \textbf{一次情報で確認できる事実}: OpenAIは2026年2月2日にCodex appを、複数のcoding agentと長時間のparallel workを管理するworkspaceとして発表した。
\item \textbf{一次情報で確認できる事実}: OpenAIは2026年2月5日にGPT-5.3-Codexを発表し、codingに加えて長時間のresearch、tool use、complex execution taskを対象とした。
\item \textbf{一次情報で確認できる事実}: GPT-5.3-Codex System Cardは2026年2月5日に公開され、OpenAIのPreparedness Frameworkに基づくcybersecurity評価をmodelへ適用している。
\item \textbf{一次情報で確認できる事実}: OpenAIは2026年2月12日にGPT-5.3-Codex-Sparkを、realtime coding interaction向けに設計したresearch-preview modelとして発表した。
\item \textbf{Vendor claim}: GPT-5.3-Codexが25\%高速、Codex-SparkがCerebras hardware上で1,000 tokens/sを超えるといった数値はvendor-reported performance claimである。
\end{itemize}
\begin{minipage}{\linewidth}
% reader-facing Technical Notes generic boundary/source tail group
{\bfseries 読む際の境界}
\begin{itemize}[leftmargin=1.5em,itemsep=0.35em]
\item \textbf{分析上の留意点}: このEvidence taskは関連する4件のfirst-party itemをまとめているため、product、base model、System Card、specialized Sparkの境界を明示する必要がある。
\item \textbf{分析上の留意点}: performanceおよびbenchmark比較はvendorによる報告である。
\item \textbf{分析上の留意点}: Codex-Sparkのthroughputに関する記述は、専用Cerebras hardwareとresearch-preview条件に依存する。
\item \textbf{分析上の留意点}: cyber capability classificationは、誇張せずSystem Cardに記載されたthreshold/precautionの境界を保持する必要がある。
\end{itemize}
\begin{samepage}
{\bfseries 一次資料}
\begingroup\sloppy
\Urlmuskip=0mu plus 2mu\relax
\begin{itemize}[leftmargin=1.5em,itemsep=0.25em]
\item {\scriptsize\url{https://openai.com/index/gpt-5-3-codex-system-card}}
\item {\scriptsize\url{https://openai.com/index/introducing-gpt-5-3-codex}}
\item {\scriptsize\url{https://openai.com/index/introducing-gpt-5-3-codex-spark}}
\item {\scriptsize\url{https://openai.com/index/introducing-the-codex-app}}
\end{itemize}
\endgroup
\end{samepage}
\end{minipage}
\endgroup
\end{technicalnote}
\medskip

\begin{technicalnote}{Claude Opus 4.6 / Sonnet 4.6}{主要資料}
\begingroup
% reader-facing Technical Notes break policy
\widowpenalty=10000
\clubpenalty=10000
\displaywidowpenalty=10000
\begin{tabularx}{\linewidth}{@{}>{\bfseries}p{0.22\linewidth}X@{}}
組織 & Anthropic \\
種別 & モデル更新 \\
時系列 & 2026-02-05 (モデル公開); 2026-02-17 (モデル公開) \\
\end{tabularx}
\smallskip
{\bfseries 一次資料から整理したtechnical points}
\begin{itemize}[leftmargin=1.5em,itemsep=0.35em]
\item \textbf{一次情報で確認できる事実}: Anthropicは2026年2月5日にClaude Opus 4.6を発表し、大規模codebaseでのcodingと継続的なagentic workを対象とする改善を示した。
\item \textbf{一次情報で確認できる事実}: Anthropicは2026年2月17日にClaude Sonnet 4.6を発表し、coding、computer use、long-context reasoning、agent planning、knowledge workの改善を示した。1M-token context windowはbetaとして説明された。
\item \textbf{Vendor claim}: Anthropicの発表に含まれるbenchmarkおよび比較性能の記述はvendor claimである。
\end{itemize}
\begin{minipage}{\linewidth}
% reader-facing Technical Notes generic boundary/source tail group
{\bfseries 読む際の境界}
\begin{itemize}[leftmargin=1.5em,itemsep=0.35em]
\item \textbf{分析上の留意点}: このEvidenceでは、関連する2件の2月model releaseを独立記事ではなく一つの系列として扱っている。
\item \textbf{分析上の留意点}: 能力およびbenchmark比較はvendorによる報告であり、本調査では独立再現していない。
\item \textbf{分析上の留意点}: availabilityと1M-contextのstatusは、発表時点のbeta/product境界を保持して解釈する必要がある。
\end{itemize}
\begin{samepage}
{\bfseries 一次資料}
\begingroup\sloppy
\Urlmuskip=0mu plus 2mu\relax
\begin{itemize}[leftmargin=1.5em,itemsep=0.25em]
\item {\scriptsize\url{https://www.anthropic.com/news}}
\end{itemize}
\endgroup
\end{samepage}
\end{minipage}
\endgroup
\end{technicalnote}
\medskip
